\documentclass[letterpaper,10pt]{article}
\usepackage[left=0.48in,right=0.48in,top=0.545in,bottom=0.415in]{geometry}
\usepackage{titlesec,enumitem,tabularx,hyperref}
\hypersetup{hidelinks,pdftitle={Nicole Jiang Resume},pdfauthor={Nicole Jiang}}
\pagestyle{empty}
\setlength{\parindent}{0pt}
\setlength{\tabcolsep}{0pt}
\raggedright
\raggedbottom
\titleformat{\section}{\scshape\large}{}{0em}{}[\titlerule]
\titlespacing*{\section}{0pt}{10pt}{5pt}
\setlist[itemize]{leftmargin=17pt,itemsep=0pt,topsep=2pt,parsep=0pt,partopsep=0pt}
\newcommand{\heading}[2]{\begin{tabular*}{\textwidth}{l@{\extracolsep{\fill}}r}\textbf{#1} & #2\end{tabular*}}
\begin{document}
\begin{center}
{\Huge\scshape Nicole Jiang}\\[8pt]
\small \href{mailto:nicolejitong.jiang@mail.utoronto.ca}{\underline{nicolejitong.jiang@mail.utoronto.ca}} $|$ \href{https://www.linkedin.com/in/nicolejitong-jiang/}{\underline{LinkedIn}} $|$ \href{https://nicolejiang.com}{\underline{nicolejiang.com}} $|$ Toronto, ON
\end{center}
\vspace{-9pt}

\section{Education}
\heading{University of Toronto}{2024 -- 2029}\\[-1pt]
{\small\textit{HBSc, Astronomy \& Physics Specialist, Statistics Minor, Philosophy Minor}}\\[2pt]
{\small\textbf{Relevant Coursework:} Computational Astrophysics, Stars and Planets, Galaxies and Cosmology, Introduction to Computer Programming, Practical Physics I, Multivariable Calculus with Proofs, Thermal Physics, Probability and Statistics I}\\[1pt]
\section{Technical Skills}
{\small\textbf{Programming:} Python (NumPy, SciPy, pandas, astropy, Matplotlib), TypeScript, SQL, Bash; Git; SSH}\\[1pt]
{\small\textbf{Data/ML:} scikit-learn; PyTorch; Transformers}\\[1pt]
{\small\textbf{Web/Tools:} React, HTML/CSS, Vite; LaTeX, Overleaf, Jupyter, VS Code, Microsoft Office, Google Workspace}\\[1pt]
\section{Research \& Technical Projects}
\vspace{5pt}\heading{Early Black-Hole Growth Constraints from JWST \href{https://github.com/nniickels/highz-accretion-atlas}{\textnormal{\underline{[GitHub]}}}}{February 2026 -- Present}\\[-1pt]
{\small\textit{Undergraduate Researcher, University of Toronto}}
\begin{itemize}\small
\item Built a provenance-tracked catalogue of 350 literature measurement records from 32 source families for accreting black holes and candidates at z $\geq$ 4, preserving uncertainties, inference methods, assumptions, and quality flags
\item Implemented vectorized NumPy growth diagnostics across seed masses, formation redshifts, Eddington ratios, spin-dependent efficiencies, and merger mass boosts
\item Quantified probabilities that required lifetime-average Eddington ratios exceed one under fixed growth histories, using 10,000 Monte Carlo draws per object with reported mass errors and mass-scale shifts of up to $\pm$0.5 dex
\item Identified 12 of 224 analysis objects requiring lifetime-average Eddington ratios above one for 100-solar-mass seeds formed at z=30, efficiency 0.1, and no merger boost; tested sensitivity to mass uncertainties and growth assumptions
\item Led catalogue development and analysis under Prof. Pratika Dayal's supervision (CITA, DAA-Dunlap)
\end{itemize}

\vspace{5pt}\heading{Galaxy Star-Formation Main Sequence Analysis with Cosmological Simulations}{May 2025}
\begin{itemize}\small
\item Queried 11,000,000+ galaxy records from public cosmological simulation databases (EAGLE, IllustrisTNG) and processed catalogues using Python, SSH sessions, and pickle serialization
\item Visualized SFR--M* trends across simulation snapshots (redshifts 0 $\leq$ z < 0.5) using scatterplots and 2D histograms
\item Compared single- vs multi-snapshot datasets to explore redshift-dependent trends in galaxy evolution
\item Drafted an AAS-style scientific paper in LaTeX using the AAS journal submission Overleaf template
\end{itemize}
\vspace{5pt}\heading{Predicting APA Site Choice from mRNA Sequences \href{https://devpost.com/software/predicting-apa-site-choice-from-mrna-sequences}{\textnormal{\underline{[Devpost]}}}}{September 2025}
\begin{itemize}\small
\item Runner-up in the Toronto Bioinformatics Hackathon as a five-member undergraduate team, placing 2/23 (\$500 prize) out of 115+ individual contestants, including many graduate students, postdocs, and industry researchers
\item Built the data preprocessing pipeline to clean and format raw RNA sequences for the DNABERT-2-117M transformer genome foundation model. Scaled the pipeline and batched inputs to train the model on a curated dataset of \textasciitilde{}800,000 sequences
\item Codified and enforced data-format standards for training; blocked non-conforming samples and implemented checks
\item Co-developed project strategy and presented results to judges; model achieved 86\% poly(A) site prediction accuracy
\end{itemize}
\vspace{5pt}\heading{Saturn and Moons Observation with HDR Imaging}{October -- December 2024}
\begin{itemize}\small
\item Captured and compiled photos of Saturn and its moons, calibrating a pixel-to-kilometre conversion to measure distances
\item Validated measurements by cross-referencing with 189 images taken over the span of two months, and comparing measurements over time to accurately evaluate Titan's orbital period
\end{itemize}
\section{Service \& Leadership}\vspace{5pt}\heading{Journal Contributor \href{https://www.rasc.ca/node/47054}{\textnormal{\underline{[Journal]}}}}{September 2025 -- Present}\\[-1pt]
{\small\textit{Royal Astronomical Society of Canada (RASC)}}
\begin{itemize}\small
\item Producing planet ephemerides for the Royal Astronomical Society of Canada's bimonthly Journal using Stellarium software
\item Curating tables of noteworthy observational astronomy events for the Journal centrespread
\item Also assisting with observatory maintenance on work trips to the E.C. Carr Astronomical Observatory
\end{itemize}
\vspace{5pt}\heading{Student Host}{February -- June 2024}\\[-1pt]
{\small\textit{Ontario Science Centre}}
\begin{itemize}\small
\item Designed and prototyped an interactive exhibition about the chemistry of art restoration, successfully pitching it to Science Centre staff to ultimately engage 100+ visitors with the final installation
\item Created and presented a puppet show to elementary school classrooms across Toronto, teaching solar-eclipse concepts
\end{itemize}

\end{document}
